\documentclass[12pt]{article}
 
\usepackage[margin=1in]{geometry} 
\usepackage{amsmath,amsfonts,amsthm,amssymb}
\usepackage{mathtools}
\usepackage{pifont}
\usepackage{graphics}
\usepackage{graphicx}
\usepackage{array}
\usepackage{listings}
\usepackage{matlab-prettifier}
\usepackage[hidelinks]{hyperref}
\usepackage{caption}
\usepackage{enumitem}

\lstdefinestyle{tex-style} {
	language=Matlab,                           % Code langugage
	basicstyle=\ttfamily,                   % Code font, Examples: \footnotesize, \ttfamily
	%keywordstyle=\color{OliveGreen},        % Keywords font ('*' = uppercase)
	commentstyle=\color{gray},              % Comments font
	numbers=none,                           % Line nums position
	numberstyle=\tiny,                      % Line-numbers fonts
	stepnumber=1,                           % Step between two line-numbers
	numbersep=5pt,                          % How far are line-numbers from code
	%backgroundcolor=\color{gray}, % Choose background color
	frame=single,                             % A frame around the code
	tabsize=2,                              % Default tab size
	captionpos=b,                           % Caption-position = bottom
	breaklines=true,                        % Automatic line breaking?
	breakatwhitespace=false,                % Automatic breaks only at whitespace?
	showspaces=false,                       % Dont make spaces visible
	showtabs=false,                         % Dont make tabls visible
	columns=flexible,                       % Column format
	morekeywords={__global__, __device__}  % CUDA specific keywords
}
\lstset{style=tex-style}
 
\newcommand{\N}{\mathbb{N}}
\newcommand{\Z}{\mathbb{Z}}
 
\newenvironment{theorem}[2][Theorem]{\begin{trivlist}
\item[\hskip \labelsep {\bfseries #1}\hskip \labelsep {\bfseries #2.}]}{\end{trivlist}}
\newenvironment{lemma}[2][Lemma]{\begin{trivlist}
\item[\hskip \labelsep {\bfseries #1}\hskip \labelsep {\bfseries #2.}]}{\end{trivlist}}
\newenvironment{exercise}[2][Exercise]{\begin{trivlist}
\item[\hskip \labelsep {\bfseries #1}\hskip \labelsep {\bfseries #2.}]}{\end{trivlist}}
\newenvironment{problem}[2][Problem]{\begin{trivlist}
\item[\hskip \labelsep {\bfseries #1}\hskip \labelsep {\bfseries #2.}]}{\end{trivlist}}
\newenvironment{question}[2][Question]{\begin{trivlist}
\item[\hskip \labelsep {\bfseries #1}\hskip \labelsep {\bfseries #2.}]}{\end{trivlist}}
\newenvironment{corollary}[2][Corollary]{\begin{trivlist}
\item[\hskip \labelsep {\bfseries #1}\hskip \labelsep {\bfseries #2.}]}{\end{trivlist}}

\newenvironment{solution}{\begin{proof}[Solution]}{\end{proof}}

%%%%%%% Table Packages %%%%%%%%%%%%%%%%%%%%%%%
\usepackage{booktabs}
\usepackage{multirow}
\usepackage{array}
\usepackage[T1]{fontenc}
\usepackage[utf8]{inputenc} 
\usepackage{tabularx,ragged2e,booktabs}
\newcolumntype{C}[1]{>{\Centering}m{#1}}
\renewcommand\tabularxcolumn[1]{C{#1}}
\newcommand{\ra}[1]{\renewcommand{\arraystretch}{#1}}
\usepackage{makecell}
\usepackage{tikz}
\usetikzlibrary{tikzmark,calc}
\usepackage{siunitx}
\usepackage[hidelinks]{hyperref}
\usepackage{amsmath}

 
\begin{document}
 
% --------------------------------------------------------------
%                         Start here
% --------------------------------------------------------------
 
\title{Homework 8}
\author{MEGN544/ROBO554 Robot Mechanics: Kinematics, Dynamics, and Control  
\\ 
Dr. George Kontoudis\\
Department of Mechanical Engineering\\
Colorado School of Mines}
\date{Fall 2026}
 
\maketitle

\begin{problem}{1}%theorem, exercise, problem, or question 
Consider a three-link manipulator with the following DH-table.
\begin{table}[!h]
\centering
\begin{tabular}{c | c | c | c | c}
\hline 
 {Link} & $a$   & $d$  & $\alpha$ & $\theta$\\  \hline
 $1$    & $0$   & $L_1$  & $\pi/2$  & $\star$ \\ 
 $2$    & $L_2$ & $0$  & $0$      & $\star$ \\ 
 $3$    & $L_3$ & $0$  & $0$      & $\star$ \\  \hline
\end{tabular}\label{tab:part_1_ii}
\end{table}

\begin{enumerate}[label=(\alph*)]
    \item Derive the forward kinematic equations symbolically for the position of the tool in the base frame as a function of $\theta_1$, $\theta_2$, and $\theta_3$.
    \item Symbolically differentiate these kinematic equations to construct a symbolic representation of the position rows of the velocity Jacobian.
    \item Using the method outlined in class, calculate $\prescript{0}{}{\boldsymbol{z}}_{k-1}$ and $\prescript{0}{}{\boldsymbol{d}}_{k-1,N}$ for all $k=1,\hdots,N$ symbolically and construct the position rows of the velocity Jacobian.
    \item For $L_1 = 1$\,m, $L_2 = 0.5$\,m, $L_3 = 0.3$\,m, $\theta_1 = 0$\,rad, $\theta_2 = \pi/4$\,rad, and $\theta_3 = \pi/4$\,rad numerically evaluate your results from (b) and (c). \textit{Hint: They should match!}
\end{enumerate}


\end{problem}

\end{document}